
%----------------------------------------------------------------------------------------
%	PACKAGES AND THEMES
%----------------------------------------------------------------------------------------

\documentclass[aspectratio=169]{beamer}
\mode<presentation> {
	
	\usetheme{Boadilla}
}
\definecolor{lmugreen}{RGB}{0, 136, 58}
\usecolortheme[named=lmugreen]{structure}
\usepackage{graphicx} % Allows including images
\usepackage{booktabs} % Allows the use of \toprule, \midrule and \bottomrule in tables
\usepackage[english]{babel}
\usepackage[utf8]{inputenc}
\usepackage{amsmath,amssymb} % math symbols
\usepackage{graphicx}
\usepackage{float}
\usepackage{hyperref} % for URLs
\newcommand{\gray}{\rowcolor[gray]{.90}}
\usepackage{verbatim}
\usepackage{textpos} % logo position
\usepackage[default]{sourcesanspro} % font type
\usepackage{cases} % math cases

% Caption
\usepackage{caption}
\DeclareCaptionFont{tiny}{\tiny}
\captionsetup{font=scriptsize,labelfont=scriptsize,justification=centering}

% ToC
\AtBeginSection[]{
	\begin{frame}[noframenumbering, plain]
	\frametitle{Outline}
	\setcounter{tocdepth}{1}
	\tableofcontents[currentsection]
\end{frame}
}

% Remove navigation 
\beamertemplatenavigationsymbolsempty

% R 
\usepackage{listings}
\lstset{
language=R,
basicstyle=\scriptsize\ttfamily,
commentstyle=\ttfamily\color{gray},
backgroundcolor=\color{white},
showspaces=false,
showstringspaces=false,
showtabs=false,
tabsize=2,
captionpos=b,
breaklines=false,
breakatwhitespace=false,
title=\lstname,
escapeinside={},
keywordstyle={},
morekeywords={}, 
belowskip = -1.2 \baselineskip,
}

%----------------------------------------------------------------------------------------
%	TITLE PAGE
%----------------------------------------------------------------------------------------

\addtobeamertemplate{frametitle}{}{%
	\begin{textblock*}{100mm}(0.88\textwidth,-0.5cm)
		\includegraphics[height=1cm,width=2cm]{../figures/lmu_logo}
\end{textblock*}}

%\includegraphics[width=0.3\textwidth]{lmu_logo}
\vspace*{-1cm}

\title{Your Title} 

\author{Your Name} 
\institute[LMU]{LMU} % Your institution as it will appear on the bottom of every slide, may be shorthand to save space
{

}
\begin{document}

{
\usebackgroundtemplate{\includegraphics[width=\paperwidth]{../figures/lmu-background.pdf}}
\begin{frame}
\begin{columns}
	\begin{column}{0.4\textwidth}
		\vspace{3cm}
		
		\textbf{\textcolor{white}{\Large Your Title}}
		\vspace{1cm}
		
		\textcolor{white}{\footnotesize Your Name \\
			\today}
	\end{column}
	\begin{column}{0.49\textwidth}
		\vspace{2cm}
		\begin{center}
			\includegraphics[width=1\textwidth]{../figures/showcase_figure.png}
		\end{center}
	\end{column}
\end{columns}
\end{frame}
}
%----------------------------------------------------------------------------------------
%	PRESENTATION SLIDES
%----------------------------------------------------------------------------------------

\section{Daten und Kontext}
\section{Fairness Definitionen}
\section{Quellen von Bias}
\section{Methoden für Fairness}

\begin{frame}{Daten und Kontext}
"begründeter Verdacht" ("reasonable suspicion") als Grund für Polizeikontrolle
"wahrscheinliche Ursache" ("probable cause") als Grund für Festnahme
mögliche Resultate:
\item Keine weiteren Maßnahmen
\item schnelle Abtastung (frisk)
\item Durchsuchung (search)
\item Festnahme (arrest)
Kritik, dass Strategie diskriminierend ggü. Schwarzen und Latinos ist
Modellieren des Entschiedungsprozesses für Festnahme mit RF
race als protected attribute (PA)
\end{frame}

\end{frame}
\begin{frame}
	\frametitle{Fairness Definitionen}
	
\end{frame}

\begin{frame}[allowframebreaks]
        \frametitle{References}
        \bibliographystyle{plain}
        \bibliography{literature.bib}
\end{frame}

\end{document}


